\documentclass[10pt,a4paper]{amsart}
\usepackage[margin=19mm]{geometry}
\usepackage{amsmath,amssymb,mathtools}
\usepackage[scr=rsfs,cal=euler]{mathalfa}
\usepackage{xfrac}
\usepackage[unicode,hidelinks,bookmarks=false]{hyperref}
\newcommand{\ie}{i.e.\ }
\newcommand{\cf}{cf.\ }
\newcommand{\resp}{respectively\ }
\newcommand{\Subsection}{\subsection}
\providecommand{\nobreakdash}{\nobreak\hbox{-}}
\setlength{\emergencystretch}{2em}
\multlinegap0pt
\usepackage{bm}
\usepackage{bbm}
\usepackage{dsfont}
\usepackage{xspace}
\usepackage{cancel}
\usepackage{mathtools}
\usepackage{amsthm}
\makeatletter
\@ifundefined{theorem}{\newtheorem{theorem}{Theorem}}{}
\makeatother
\title[Beatty Sequences for a Quadratic Irrational: Decidability an]{Beatty Sequences for a Quadratic Irrational: Decidability and Applications}
\author{Luke Schaeffer, Jeffrey Shallit, Stefan Zorcic}
\date{Source: arXiv:2402.08331v3 (CC BY 4.0)}
\begin{document}
\maketitle

\noindent\emph{Source and licence.} Beatty Sequences for a Quadratic Irrational: Decidability and Applications. Version arXiv:2402.08331v3 (CC BY 4.0), distributed under the Creative Commons Attribution 4.0 International licence, as recorded in the arXiv licence metadata for this exact version and shown on the abstract page https://arxiv.org/abs/2402.08331v3. Full rights notice: licenses/arxiv-2402.08331-CC-BY-4.0.txt.\par\smallskip

\noindent\textit{Editorial scope.} Counted unit: the Theorem whose source label is \texttt{None} in the article (source0.tex lines 776--796, source label \texttt{None}), delivered together with the article's own complete original proof of it (source0.tex lines 798--829, one \texttt{proof} environment of 32 lines). No mathematics is written, repaired or re-derived: statement and proof are the article's own text line for line. Required context retained uncounted: varphisynch (lines 696--699). Every definition, display and label of those lines is retained unchanged. Omitted from this package and not redistributed: 1--695, 700--775, 797--797, 830--1880. Nothing inside a retained excerpt is omitted or altered: every retained body is delivered exactly as the article's lines are written, so no omission receipt of any kind accompanies this package. Citations: the retained excerpts carry no citation command at all, so no bibliography entry of the article is transcribed and no reference is redistributed. Reference adaptation: none was needed and none was made; every reference inside the delivered unit resolves inside it, so no unresolved reference or citation is produced. Printed slips kept literal: no printed word, symbol, hypothesis or number of the counted statement or of its proof was altered. Layout and class: the article's own class is \texttt{article} with packages including color, amssymb, amsmath, amsthm, amsfonts, amscd, graphicx, hyperref, fullpage, float, graphics, latexsym, epsf; the wrapper is the lane's pinned \texttt{amsart} editorial wrapper at 10pt on A4, which restates the theorem-like environments the retained text needs and adds the article's own macro definitions that the retained text needs (none), with extra wrapper packages (bm, bbm, dsfont, xspace, cancel, mathtools). The delivered numbering is the unit's own. Assets: no retained span contains a figure environment, a graphics-inclusion command, a PDF-page inclusion, a table or a TikZ picture; no image file is copied into this package, the package declares no asset dependency and no image file is redistributed with it. Licence: version arXiv:2402.08331v3 of this article is distributed under the Creative Commons Attribution 4.0 International licence, as recorded in the arXiv licence metadata for this exact version and shown on the abstract page https://arxiv.org/abs/2402.08331v3. A full rights notice is delivered at licenses/arxiv-2402.08331-CC-BY-4.0.txt. Characters: every retained excerpt is pure ASCII, so the delivered engine, XeLaTeX with its default Latin Modern text font, reports no missing character.
\begin{theorem}
The sequence $\bf c$ is Fibonacci synchronized.
\label{varphisynch}
\end{theorem}

\begin{theorem}
Let $r = 2$, $s = (\sqrt{5}-1)/2$, and $t = (1+\sqrt{5})/2$.
Define
\begin{align*}
a(n) &= n + \lfloor ns/r \rfloor + \lfloor nt/r \rfloor \\
b(n) &= n + \lfloor nr/s \rfloor + \lfloor nt/s \rfloor \\
c(n) &= n + \lfloor nr/t \rfloor + \lfloor ns/t \rfloor .
\end{align*}
Then 
\begin{itemize}
\item[(a)]
there exists $m$ such that $m = a(n)$ iff (i) holds
for $r= n+1$ but (ii) does not for $r = n+1$;
\item[(b)]
there exists $m$ such that $m = b(n)$ iff neither
(i) nor (ii) holds for $r = n+1$;
\item[(c)]
there exists $m$ such that $m = c(n)$ iff both
(i) and (ii) hold for $r = n+1$.
\end{itemize}
\end{theorem}

\begin{proof}
We implement the formulas above in \texttt{Walnut}.   Here {\tt F} is {\tt Walnut}'s abbreviation for the Fibonacci word, indexed starting at position $0$,
and {\tt phin} is the synchronized automaton computing
$\lfloor \varphi n \rfloor$ for $\varphi = (1+\sqrt{5})/2$, the golden ratio,
given in the proof of Theorem~\ref{varphisynch}.
\begin{verbatim}
def three0 "?msd_fib Ei (F[i]=@0) & (F[i+n]=@0) & (F[i+2*n]=@0)":
def three1 "?msd_fib Ei (F[i]=@1) & (F[i+n]=@1) & (F[i+2*n]=@1)":

def fibsr "?msd_fib Er $phin(n,r) & s=(r-n)/2":
def fibtr "?msd_fib Er $phin(n,r) & s=r/2":
def fibrs "?msd_fib $phin(2*n,s)":
def fibts "?msd_fib Er $phin(n,r) & s=r+n":
def fibrt "?msd_fib Er $phin(2*n,r) & s=r-2*n":
def fibst "?msd_fib Er $phin(n,r) & s=2*n-(r+1)":

def a189377 "?msd_fib En,x,y $fibsr(n,x) & $fibtr(n,y) & z=n+x+y":
def a189378 "?msd_fib En,x,y $fibrs(n,x) & $fibts(n,y) & z=n+x+y":
def a189379 "?msd_fib En,x,y $fibrt(n,x) & $fibst(n,y) & z=n+x+y":

def reble1 "?msd_fib $three0(n) & ~$three1(n)":
def reble2 "?msd_fib ~$three0(n) & ~$three1(n)":
def reble3 "?msd_fib $three0(n) & $three1(n)":

eval rebleconj1 "?msd_fib An (n>=1) => ($reble1(n+1) <=> $a189377(n))":
eval rebleconj2 "?msd_fib An (n>=1) => ($reble2(n+1) <=> $a189378(n))":
eval rebleconj3 "?msd_fib An (n>=1) => ($reble3(n+1) <=> $a189379(n))":
\end{verbatim}
When we run them, we get 
\texttt{TRUE}
as the result.   The entire computation takes less than a second.
\end{proof}

\end{document}
